% Ready-to-paste insertions for the uploaded TCSSOP manuscript.
% Each block shows WHERE to cite one Computers & Industrial Engineering paper.

%--------------------------------------------------------------------------
% 1) Afzalirad & Rezaeian (2016)
% Place: Introduction -- OR/scheduling difficulty paragraph
%--------------------------------------------------------------------------
% Original idea: eligibility + tardiness heuristics.
% Suggested sentence ending:
%
% ``...unrelated parallel-machine scheduling with machine eligibility and
% tardiness objectives is known to be computationally challenging and often
% requires specialized heuristic methods
% [Ulaga et al., 2022, Maecker et al., 2023, Afzalirad and Rezaeian, 2016].''
%
% Bibliography entry:
% M. Afzalirad and J. Rezaeian. Resource-constrained unrelated parallel
% machine scheduling problem with sequence dependent setup times, precedence
% constraints and machine eligibility restrictions. Computers \& Industrial
% Engineering, 98:40--52, 2016. doi: 10.1016/j.cie.2016.05.020.

%--------------------------------------------------------------------------
% 2) Wang et al. (2017)
% Place: Literature Review -- multi-project / priority-rule stream
%--------------------------------------------------------------------------
% Suggested added sentence after RCMPSP discussion:
%
% ``Priority-rule studies further show that due-date-oriented dispatching
% rules, including earliest-due-date variants, remain competitive baselines
% for multi-project settings under resource contention
% [Wang et al., 2017].''
%
% Also usable in Section 4.4 (EDD scheduler motivation).
%
% Bibliography entry:
% Y. Wang, Z. He, L.-P. Kerkhove, and M. Vanhoucke. On the performance of
% priority rules for the stochastic resource constrained multi-project
% scheduling problem. Computers \& Industrial Engineering, 114:223--234,
% 2017. doi: 10.1016/j.cie.2017.10.021.

%--------------------------------------------------------------------------
% 3) Hulett et al. (2017)
% Place: Literature Review -- chamber / laboratory analogy
%--------------------------------------------------------------------------
% Suggested added sentence after eligibility discussion:
%
% ``Closely related laboratory and burn-in settings further motivate modeling
% shared environmental resources as non-identical parallel batch processors
% with tardiness objectives [Hulett et al., 2017].''
%
% Bibliography entry:
% M. Hulett, P. Damodaran, and M. Amouie. Scheduling non-identical parallel
% batch processing machines to minimize total weighted tardiness using
% particle swarm optimization. Computers \& Industrial Engineering,
% 113:425--436, 2017. doi: 10.1016/j.cie.2017.09.037.

%--------------------------------------------------------------------------
% 4) Perez-Gonzalez et al. (2019)
% Place: Section 4.4 -- Scheduling procedure (constructive EDD)
%--------------------------------------------------------------------------
% Suggested added sentence after introducing the EDD ready-set rule:
%
% ``Constructive heuristics of this type are particularly attractive for
% unrelated parallel-machine problems with eligibility and due-date pressure,
% where repeated decoding must remain computationally light
% [Perez-Gonzalez et al., 2019].''
%
% Bibliography entry:
% P. Perez-Gonzalez, V. Fernandez-Viagas, M. Zamora Garcia, and J. M. Framinan.
% Constructive heuristics for the unrelated parallel machines scheduling
% problem with machine eligibility and setup times. Computers \& Industrial
% Engineering, 131:131--145, 2019. doi: 10.1016/j.cie.2019.03.034.

%--------------------------------------------------------------------------
% 5) Liao & Cheng (2007)
% Place: Literature Review VNS paragraph AND/OR Section 4.5 VNS method
%--------------------------------------------------------------------------
% Suggested Literature Review ending:
%
% ``...combine local improvement with systematic diversification
% [Mejia and Yuraszeck, 2020, Liao and Cheng, 2007].''
%
% Suggested Method sentence:
%
% ``Following the classical VNS principle of changing neighborhood intensity
% after local stagnation [Liao and Cheng, 2007], the method uses a fixed
% fallback ratio sequence...''
%
% Bibliography entry:
% C.-J. Liao and C.-C. Cheng. A variable neighborhood search for minimizing
% single machine weighted earliness and tardiness with common due date.
% Computers \& Industrial Engineering, 52(4):404--413, 2007.
% doi: 10.1016/j.cie.2007.01.004.
